\section{The result and its scope}\label{sqt:sec:results}
For $n\geq1$, define
\[
a_n=\#\left\{(m_{ij})\in\mathbb Z_{\geq0}^{n\times n}:
 \sum_jm_{ij}=n\text{ for all }i,\quad
 \sum_im_{ij}=n\text{ for all }j\right\}.
\]
Rows and columns are labelled. Thus matrices related by a permutation of rows or columns are generally distinct. The first values are $1,3,55,10147,22069251$ for $n=1,\ldots,5$; with $a_0=1$, this is OEIS A110058 \cite{oeis}. Throughout, all logarithms are natural, and asymptotic assertions concern integer $n\to\infty$ unless a real variable is specified. Set
\begin{equation}\label{sqt:eq:leading}
G_n=\frac{4^{n^2}}{(4\pi)^{n-1/2}n^{n-1}},\qquad L_n=e^{1/4}G_n.
\end{equation}
The leading equivalent $a_n\sim L_n$ is already contained in Canfield--McKay \cite{CM}; it is also covered by the smooth-margin theory of Barvinok--Hartigan \cite{BH}. The purpose here is to justify every fixed algebraic order in this square density-one specialization, evaluate two corrections, and control inversion.

\begin{theorem}\label{sqt:thm:main}
There are rational numbers $C_j$, $j\geq1$, such that for every fixed integer $K\geq0$,
\begin{equation}\label{sqt:eq:main}
\log\frac{a_n}{L_n}=\sum_{j=1}^{K}\frac{C_j}{n^j}+O_K(n^{-K-1}).
\end{equation}
Each $C_j$ is given by a finite rational combinatorial algorithm. In particular,
\begin{equation}\label{sqt:eq:two}
C_1=-\frac32,\qquad C_2=\frac{223}{32}.
\end{equation}
Consequently,
\begin{align}
\log\frac{a_n}{L_n}&=-\frac{3}{2n}+\frac{223}{32n^2}+O(n^{-3}),\label{sqt:eq:logtwo}\\
\frac{a_n}{L_n}&=1-\frac{3}{2n}+\frac{259}{32n^2}+O(n^{-3}).\label{sqt:eq:relativetwo}
\end{align}
\end{theorem}

The effective assertion is about individual coefficients. It does not assert explicit values of the implicit remainder constants or starting indices. The theorem is a Poincar\'e expansion at every fixed order, with $K$ fixed before $n\to\infty$. It does not assert convergence of the infinite series, uniformity in growing $K$, optimal truncation, or a description of exponentially small contributions.

The argument has three distinct parts. First, an exact Fourier integral is localized by the integer-table theorem of Canfield--McKay. A Gaussian gauge penalty and a linear whitening map then convert the localized problem into an expectation on a product of truncated Gaussian intervals. Second, Isaev's theorem for bounded \emph{complex} functions \cite{Isaev} gives a sufficiently accurate finite cumulant expansion. Third, connected Wick contractions provide exact Laurent polynomials and a degree bound that leaves only finitely many degree multisets at any required order. The complex cumulant method is imported; the checks needed to apply it to these geometric tables are supplied in full.

For later use, define the inverse counting function
\begin{equation}\label{sqt:eq:inversecount}
\nu(Y)=\min\{n\geq1:a_n\geq Y\},\qquad Y>1,
\end{equation}
and, for fixed $K$, the real approximation
\begin{equation}\label{sqt:eq:FK}
F_K(x)=x^2\log4-(x-\tfrac12)\log(4\pi)-(x-1)\log x
       +\frac14+\sum_{j=1}^{K}C_jx^{-j}.
\end{equation}
There is a unique sufficiently large solution $t_K=t_K(Y)$ of $F_K(t_K)=\log Y$. We prove in Section~\ref{sqt:sec:inverse} that some $D_K>0$ gives, for all sufficiently large $Y$,
\begin{equation}\label{sqt:eq:inverseintro}
\ceil{t_K-D_Kt_K^{-K-2}}\leq\nu(Y)\leq
\ceil{t_K+D_Kt_K^{-K-2}}.
\end{equation}
The two ceilings matter near an integer threshold. An unqualified rule $\nu(Y)=\lceil t_K\rceil$ is not a consequence of the theorem.

\begin{remark}[{[write]} The OEIS entry, 6~October 2026]\label{sqt:rem:oeis}
The live entry A110058 \cite{oeis} (revision~\#30 of 30~May 2026, read
6~October 2026) is named ``Number of nonnegative integer matrices of order n
for which all row and column sums equal n.'' Its data run through $a_{13}$,
its one link is to the Combinatorica version of \cite{CM}, and its only
formula, the only asymptotic statement in the entry, reads verbatim
\begin{center}\small
\texttt{log a(n) = 2(log 2)*n\^{}2 - n*(log n) - n*(log 4*Pi) + (log n) + O(1). - \_Igor Pak\_, May 15 2019}
\end{center}
(with \texttt{log 4*Pi} read as $\log(4\pi)$). Since
$\log L_n=n^2\log4-(n-\tfrac12)\log(4\pi)-(n-1)\log n+\tfrac14$, Pak's main
terms are exactly $\log L_n-\tfrac12\log(4\pi)-\tfrac14$. So the formula is
consistent with $a_n\sim L_n$, which is Canfield--McKay's, and the latter
identifies Pak's bounded term as $\tfrac12\log(4\pi)+\tfrac14+o(1)
=1.51551\ldots+o(1)$; Theorem~\ref{sqt:thm:main} expands the $o(1)$ to every
fixed order. Nothing in the entry is corrected, and nothing was submitted to
the OEIS. The ``b-file'' cited by the source is synthesized by the OEIS from
the entry's data (its header reads ``A110058 (b-file synthesized from
sequence entry)'', $0\le n\le13$); the counts at $n=6,8,10,13$ in
Table~\ref{sqt:tab:diagnostics} are taken from it.
\end{remark}

\subsection{Provenance, sources and reading conventions}\label{sqt:sec:provenance}
\begin{writenote}[Added 6 October 2026, batch~108]
\emph{Provenance.} This report is Report~236 of the session bundle that
ProveIt's \file{docs/incoming} intake processed as batch~108: the archive
\file{Report236.zip} (669{,}911 bytes, SHA-256
\file{1b73cd6bd55d...5b132059}, 35 files in a wrapper directory
\file{Report236/}), arrival commit \file{60f54ea06}, placed in
\file{602e5bd0f}. The text is the delivered \file{article.tex} with its
twelve files \file{sections/*.tex} (980 lines in all, dated 5~October 2026),
shipped under the same names. The delivered 24-page PDF and the checksum
manifest \file{MANIFEST.sha256} (34 entries, verified at placement) are not
shipped, and the delivered README is replaced by the report README; all
three are retrievable from the arrival commit. The package records no
ProveIt commit and cites nothing in the repository. The phrase ``private
review'' occurs in it only in sentences saying that such material is not
included (Section~\ref{sqt:sec:computation}, \file{SOURCES.md} and the
delivered README); the PDF author field is ``Report 236''. It is a
single-source report, so no merge choice was made. The write prefixed the
97 delivered labels with \file{sqt:} and updated the 87 references to them;
it added this subsection, the notes marked [write],
Remarks~\ref{sqt:rem:oeis}, \ref{sqt:rem:uniformity},
\ref{sqt:rem:cmregime} and~\ref{sqt:rem:transseries},
Lemma~\ref{sqt:lem:sequential} and Section~\ref{sqt:sec:further}, and it
narrowed three sentences on the correction-factor conjecture (the abstract,
the title of Section~\ref{sqt:sec:conjecture} and the sentence after
Corollary~\ref{sqt:cor:delta}; first wordings kept in the notes). No
theorem, proof or number of the source was changed: the added statements are
the last of their sections, the added section comes after the last numbered
one, and the added displays are unnumbered, so every theorem, section,
table and equation keeps its delivered number.

\emph{The sources, as the write read them} (6~October 2026).
\begin{itemize}\raggedright
\item OEIS A110058, revision~\#30, and its synthesized b-file
 (Remark~\ref{sqt:rem:oeis}).
\item Canfield--McKay \cite{CM}, arXiv:math/0703600v2: pages 1--6 and
 15--20 (Theorem~1 with condition~(1.1), Conjecture~1 with~(1.4) and their
 record of proved cases, the convention on asymptotic constants, the
 integral~(2.2)--(2.3), the region~(3.1), Theorem~3, Lemmas~3 and~4 and
 the proof of Theorem~3). Checked against the source: at $m=n$, $s=t=\lambda
 n$ their integrand $F$ of~(2.3) is $\mathcal F_\lambda$ of
 Section~\ref{sqt:sec:density}, since
 $e^{-\ii s\sum_j\theta_j-\ii t\sum_k\phi_k}=\prod_{j,k}e^{-\ii\lambda(\theta_j+\phi_k)}$;
 their $I_0$ of~(4.1) is $I_{0,\lambda}$ of~\eqref{sqt:eq:densityminor}
 (and $I_0$ of~\eqref{sqt:eq:minor} at $\lambda=1$); their region~(3.1) is
 the region described in Sections~\ref{sqt:sec:imports}
 and~\ref{sqt:sec:density}, with their $\varepsilon$ equal to the source's
 $\eta$; the left side of~(1.1) is $3$ at $\lambda=1$ and
 $\tfrac23(4+1/u)$ in general; and their leading formula~(1.3) at $m=n$ is
 $\mathcal G(n,\lambda)e^{B_0(u)}$, so $B_0$ and $\mathcal G$ (and $L_n$ at
 $\lambda=1$) are theirs. Their Theorem~3 gives $O(e^{-n^\varepsilon})I_0$;
 Proposition~\ref{sqt:prop:minor} states the weaker $O(e^{-cn^\eta})I_0$.
 Not read: their Sections~3 and~5 (the central evaluation and the proof of
 Theorem~1), which the source does not use.
\item Isaev \cite{Isaev}, arXiv:2508.16952v2 (HTML version): the statement
 of Theorem~1.2 and the definitions of $\Delta_V$, $\overline\Delta_V$ and
 $S_k$. Proposition~\ref{sqt:prop:complex} is that theorem with $n=N$ and
 $S_k$ replaced by the larger $D_k$, since
 $\overline\Delta_V(f,X)\le\Delta_V(f)$.
\item The transseries volume
 \path{Analysis/Transseries/docs/series-and-transseries/Transseries_And_Inversion/transseries_and_inversion.tex}:
 \file{p0:def:core}, \file{p0:thm:lambert-core},
 \file{p0:prop:factorial-core}, \file{p0:thm:core-reversion},
 \file{p0:def:three-inverses}, \file{p0:thm:staircase} and
 \file{plt:thm:lw-template}, read for Remark~\ref{sqt:rem:transseries}; and,
 at the independent check of 6~October 2026,
 \file{plt:def:lw-monomial-log-datum}.
\item Not read by the write: Barvinok--Hartigan \cite{BH} (beyond its arXiv
 title page), Greenhill--McKay \cite{GM}, Isaev--McKay \cite{IM,IM2018},
 Isaev--McKay--Zhang \cite{IMZ}, Zipunnikov--Booth--Yoshida \cite{ZBY}, Aw
 \cite{Aw} and Isaev--Makai--McKay \cite{IMM}. What is said about them
 (Barvinok--Hartigan's Theorem~1.3 and the comparison~\eqref{sqt:eq:BHcompare},
 the hypotheses of Isaev--McKay, Greenhill--McKay's Corollaries~4.1
 and~4.2) is the source's. Canfield--McKay's own record (their case~(b),
 quoted in Remark~\ref{sqt:rem:cmregime}) independently states that the
 conjecture holds eventually in Greenhill--McKay's sparse range.
\end{itemize}

\emph{What was checked at intake.} The batch-108 dossier read the
manuscript, checked the conjecture and condition~(1.1) against the
Canfield--McKay text, and reran the package on a copy (Windows, Python
3.14.4, SymPy 1.14.0): \file{build.py --verify-only} verified 34 files;
\file{connected_wick.py}, \file{verify_coefficients.py} and
\file{count_tables.py} reproduced the delivered receipts after line-ending
conversion; \file{diagnostics.py} reproduced
Table~\ref{sqt:tab:diagnostics}; \file{guard_tests.py} stops on Windows
at its own path guard; the PDF and ZIP rebuilds were not run. It
computed $\Delta(n,n;n,n)$ from exact counts and found $1.070,\ldots,0.565$
for $n=4,5,6,8,10,13$, and Richardson extrapolation of $D_1,D_2$ over
$6\le n\le13$ gave about $-1.47$ and $6.5$ against $-3/2$ and $223/32$. The
write rechecked by hand the coefficient identity
$P_2(1/u)-g_2(u)-\tfrac56=\tfrac{67}{12}+\tfrac5{6u}$, the expansion of
$(n-1)\log(1+1/n)$, the two normalization identities of
Section~\ref{sqt:sec:density}, the value $3$ of condition~(1.1) and the
identifications listed above; it computed $\Delta(n,n;n,n)$ for all
$1\le n\le13$ (Remark~\ref{sqt:rem:cmregime}) and checked in SymPy that the
four reversion coefficients of~\eqref{sqt:eq:reversion} solve the
fixed-point equation of Remark~\ref{sqt:rem:transseries}. Finite
computations are observations, not certificates.

\emph{Relation to the repository.} Before batch~108 no file of the
repository mentioned A110058 or Canfield and McKay's conjecture. No
statement of this report is formalized in Lean or Rocq, and its place in the
collection confers no formal status. The neighbouring reports are listed in
the report README; \file{a380592-tied-football-seasons} applies the same
Isaev theorem to another integral and receives a reciprocal note.
\end{writenote}
\begin{writenote}[What is not claimed, collected from the source]
No numerical value of a remainder constant or a starting index is asserted:
Theorem~\ref{sqt:thm:main} is a Poincar\'e expansion at each fixed order $K$,
with no convergence of the series, no uniformity in growing $K$, no optimal
truncation and no exponentially small terms; no nonperturbative sectors,
resurgence or Stokes behaviour are established. The leading equivalent is
Canfield--McKay's and is also within Barvinok--Hartigan's theory; Isaev's
complex cumulant theorem and the connected-Wick toolkit are imported, and
no new general saddle-point or cumulant method is claimed. The literature
review was bounded: absence of the coefficient pair from the sources
inspected ``is not a worldwide priority certificate''; Greenhill--McKay's
sparse result is the antecedent of the first coefficient, and
equation~\eqref{sqt:eq:BHcompare} is a specialization of Barvinok--Hartigan's
formula, not a correction. The inverse constants are existential:
$\lceil t_K\rceil$ is not a valid rule, computing a root does not compute
$B_K$, and the Newton statement is about exact-real iteration. The density
extension is not uniform as the density tends to $0$ or~$\infty$; at a fixed
irrational density there are no exact margins. The corollary on the
correction factor has no explicit starting index and does not settle finite
sizes, nonsquare tables or degenerating densities. No monotonicity of the
diagnostics or sign of omitted terms is claimed. The package's computations
verify finite identities and small counts in a fixed validated range (cost at
most three); a receipt certifies no remainder, localization, finite-size
inequality, inverse ceiling or priority; byte-identical PDF reproduction is
toolchain-specific; the manifest is ``not a digital signature''; and the
reproducibility controls are ``not a sandbox''. The write adds: the
uniformity in density of the imported localization is established in this
text only in the weaker form of Lemma~\ref{sqt:lem:sequential}
(Question~\ref{sqt:q:uniform}).
\end{writenote}
\medskip\noindent\begin{minipage}{\textwidth}
\emph{Reading conventions} (letters the manuscript reuses, with the tempting
false readings; no symbol of the source was renamed; symbols of the
transseries volume carry the subscript ``vol'' in
Remark~\ref{sqt:rem:transseries}).\par
\smallskip
\begingroup\small\renewcommand{\arraystretch}{1.1}
\noindent\begin{tabular}{@{}>{\raggedright\arraybackslash}p{0.85in}>{\raggedright\arraybackslash}p{3.15in}>{\raggedright\arraybackslash}p{2.15in}@{}}
\toprule
Symbol & Meaning here (where) & Not to be read as\\
\midrule
$\Delta$ & Isaev's mixed difference $\Delta_V(f)$ (Section~\ref{sqt:sec:imports}); Canfield--McKay's correction $\Delta(m,s;n,t)$ (Section~\ref{sqt:sec:conjecture}) & each other\\
$D_k$, $D_K$, $D_1$, $D_2$ & Isaev's sums $D_k(f)$; the constant of Theorem~\ref{sqt:thm:inverse}; the diagnostics of Section~\ref{sqt:sec:computation} & each other\\
$G_n$, $G_{\rm Good}$, $\mathcal G$ & Gaussian prefactor~\eqref{sqt:eq:leading}; Canfield--McKay's (Good's) binomial estimate; density prefactor~\eqref{sqt:eq:densityG} & each other\\
$A$, $A_\lambda$ & augmented matrix~\eqref{sqt:eq:augmented}; the scalar $u/2$, Canfield--McKay's $A$ & each other\\
$T$, $T(n,s)$ & whitening map~\eqref{sqt:eq:whitening}; the table count of Section~\ref{sqt:sec:density} & each other\\
$h$ & local interaction~\eqref{sqt:eq:h}; the gauge coordinate $x_{2n}$ (Section~\ref{sqt:sec:gauge}) & $h_{\rm vol}$ of \file{p0:thm:core-reversion}\\
$m$, $L$ & cumulant order and Taylor cutoff~\eqref{sqt:eq:cutofforders} & Canfield--McKay's row count $m$ ($=n$ here); $L_n$ of~\eqref{sqt:eq:leading}\\
$\varepsilon$, $\eta$ & cube exponent; localization exponent $\eta=\varepsilon/2$ & Canfield--McKay's $\varepsilon$, which is the source's $\eta$\\
$a$, $b$ & in Section~\ref{sqt:sec:imports}: the constants of condition~(1.1); in Section~\ref{sqt:sec:inverse}: $\log4$, $\log(4\pi)$ & each other; $a_n$\\
$r$ & number of cumulant arguments; $\sqrt{\log Y/a}$ (Section~\ref{sqt:sec:inverse}); exceptional count (Section~\ref{sqt:sec:density}) & each other\\
$u$ & $\lambda(1+\lambda)$; the gauge parameter in the proof of Lemma~\ref{sqt:lem:transfer} & $u_{\rm vol}$ of \file{p0:thm:core-reversion}\\
$d,e,f,g$ & reversion coefficients~\eqref{sqt:eq:reversiond}--\eqref{sqt:eq:reversiong} & degrees $d_a$; $f=R\circ T$\\
$P_j$, $\Part$ & density polynomials; set partitions & each other; $P_a(A)$ power sums\\
$\nu(Y)$ & $\min\{n\ge1:a_n\ge Y\}$ & it \emph{is} the staircase $N_\ast$ of \file{p0:def:three-inverses} (Remark~\ref{sqt:rem:transseries})\\
$\kappa$ & cumulants $\kappa_r$; geometric cumulants $\kappa_d(\lambda)$ & the $\kappa$ of \file{p0:prop:factorial-core}\\
\bottomrule
\end{tabular}
\endgroup
\end{minipage}\par
